%------------------------------------------------------------------------------%
\section{Divisão de polinômios}

Dividir um polinômio $f(x)$ de grau $n$ por um polinômio $g(x)\neq 0$ de grau
menor que $n$ é encontrar dois polinômios $q(x)$ e $r(x)$ tais que
\[
  f(x) = g(x) q(x) + r(x)
\]
O grau de $r(x)$ é menor que o grau de $g(x)$.

%------------------------------------------------------------------------------%
\begin{example}
  Dividir $f(x)=6x^3-2x^2+x+3$ por $g(x)=x^2-x+1$.

  \solution
  \[
    \begin{array}{r|r}
      \dropsign{-} 6x^3 - 2x^2 + \phantom{6}x + 3 & x^2 - \phantom{6}x + 1 \\ \cline{2-2}
      6x^3 - 6x^2 + 6x \phantom{{}+3}             & 6x + 4                 \\ \cline{1-1}
                                                  &  \\[\dimexpr-\normalbaselineskip+\jot]
      \dropsign{-} 4x^2 - 5x + 3                  &  \\
      4x^2 - 4x + 4                               &  \\ \cline{1-1}
                                                  &  \\[\dimexpr-\normalbaselineskip+\jot]
      -  x - 1                                    &
    \end{array}
  \]
  Logo,
  \[
    6x^3 - 2x^2 + x+3 = (x^2-x+1)(6x+4) - x-1.
  \]
\end{example}
%------------------------------------------------------------------------------%

